\begin{abstract}
Background work on desktops (sync clients, indexers, updates, builds) degrades the tail latency of interactive applications. Windows gives unprivileged software no way to measure that harm or to tell which resource causes it. We present SmartSwap, a user-mode system that (i)~estimates the share of wall time an interactive task loses to interference with a \emph{Windows Stall Index} (WSI) measured by a low-duty canary, (ii)~attributes the stall to CPU, memory bandwidth, or storage with a specificity-ordered differential rule, and (iii)~applies a least-cost ladder of documented, reversible controls (scheduler hints, then an AIMD-governed Job Object CPU-rate cap). Every action is verified twice: by kernel read-back that it took effect, and by the canary that it helped. In a 340-trial randomised block study on a hybrid-core laptop, WSI tracked the tail latency of a separate, un-instrumented foreground with Spearman $\rho=0.96$. It detected interference with ROC AUC 0.965--1.000 and a 1.2\,\% false-alarm rate, and it named the injected resource in 86--99.6\,\% of stalled periods. The controller cut foreground p95 by 87\,\% under CPU and mixed contention, and was the only policy to help under storage contention (0.55$\times$ the p95 of a CPU-cap policy). Removing the ladder made p95 up to 4.8$\times$ worse. All 1\,575 actions and 322 reverts were verified. Static priority demotion proved an equally effective baseline for latency, and we report where adaptive control over-throttled. We also dissect how our own first evaluation produced convincing but invalid results.
\end{abstract}
